\documentclass[14pt, aspectratio=1610]{beamer}
%\documentclass[12pt,aspectratio=1610,handout]{beamer}
\usepackage{csu-theme}
\usepackage{binary-trees}
\usepackage{arrays}
\usepackage{tabularx}

\title{\Huge{} Priority Queues}
\subtitle{CSCI 315: \textit{Data Structure Analysis}}
\author{Sean T. Hayes, Ph.D.}
\institute{Charleston Southern University}

%% Comment out date to use default (today)
% \date{April 1, 2025}
\subject{Software Programming}
%\keywords{}

\setlength{\parskip}{0.5em}

\begin{document}

\begin{frame}
  \titlepage{}
\end{frame}

\section*{Intro}

\begin{frame}{Introduction}
  \begin{itemize}[<+->]
  \item
    A stack is first in, last out.
  \item
    A queue is first in, first out.
  \item
    A priority queue is a \textit{highest-priority first-out} container.
  \end{itemize}
\end{frame}

\begin{frame}{Introduction}
  \begin{itemize}[<+->]
  \item
    The ``highest-priority'' element is removed first.\\
    Priority is determined from the elements' values,\\ usually largest to smallest or vice versa.
  \item
    The definition of ``largest'' is up to the programmer (for example, you might define it by implementing a \href{https://en.cppreference.com/w/cpp/named_req/Compare}{Compare} function object.)
  \item
    If there are several ``largest'' elements, the implementer must decide which to remove first.
    \begin{itemize}
    \item
      Remove any ``largest'' element (don't care which).
    \item
      Remove the first one added.
    \end{itemize}
  \end{itemize}
\end{frame}

\begin{frame}{Popular Priority Queue Uses}
  \begin{itemize}[<+->]
  \item
    \emph{Task scheduling} in operating systems (highest-priority processes run first).
  \item
    \emph{Dijkstra's algorithm} for shortest paths (select the next closest node).
  \item
    \emph{A* search} and other \emph{graph traversal} techniques (best candidate expansion).
  \item
    \emph{Event-driven simulation} (next event by earliest time).
  \item
    \emph{Huffman coding} and \emph{compression} (combine nodes by lowest frequency).
  \end{itemize}
\end{frame}

% \setbeamercovered{transparent}
\begin{frame}{A Priority Queue ADT}
Here is one possible ADT\@:
\renewcommand{\arraystretch}{1.25}
\onslide<2->{%
\begin{tabularx}{1\textwidth}{>{\columncolor{codeBackground}\color{codeVariable}\ttfamily}l >{\raggedright\arraybackslash} X}
  {PriorityQueue{\color{white}()}} & a constructor \onslide<3->{\\
  {{\color{codeKeyword}void} add{\color{white}(}{\color{codeKeyword}const} Type {\color{white}\&}val{\color{white})}} & inserts \texttt{val} into the queue} \onslide<4->{\\
  {Type{\color{white}\&} pop{\color{white}()}} & removes and returns the highest-priority element} \onslide<5->{\\
  \texttt{{\color{codeKeyword}const} Type{\color{white}\&} top{\color{white}()} {\color{codeKeyword}const}} & returns (but does not remove) the highest-priority element} \onslide<6->{\\
  {{\color{codeKeyword}bool} isEmpty{\color{white}()} {\color{codeKeyword}const}} & returns true IFF empty} \onslide<7->{\\
  {{\color{codeKeyword}unsigned int} size{\color{white}()} {\color{codeKeyword}const}} & returns the element count} \onslide<8->{\\
  {{\color{codeKeyword}void} clear{\color{white}()}} & discards all elements}
\end{tabularx}%
}
\end{frame}

\setbeamercovered{transparent}


\setbeamercovered{}

\section{Performance}

\begin{frame}{Evaluating Implementations}
\begin{itemize}[<+->]
\item
  When choosing a data structure, consider \\\emph{usage patterns}.
  \begin{itemize}
  \item
    If data is loaded once and searched thousands of\\
    times, make searching fast by sorting the array.
  \item
    If we plan to perform do only a few searches on a huge\\ data set, we probably \emph{don't} want to spend time sorting.
  \end{itemize}
\item
  For almost all uses of a queue (including a priority queue), we eventually remove everything.
\item
  Hence, when we analyze a priority queue, neither `add' nor ``remove'' is more important --- we must look at the timing for ``add + remove''.

\end{itemize}
\end{frame}

\begin{frame}{Array Implementations: Unsorted}
  A priority queue could be implemented as an \emph{unsorted}\\ array (with a count of elements).
  \pause{}
  \begin{itemize}[<+->]
    \item
      Adding an element would take \(O(1)\) time. (Why?)
    \item
      Removing an element would take \(O(n)\) time. (Why?)
    \item
      Hence, adding and removing an element takes \(O(n)\) time.
    \item
      This is an inefficient representation.
  \end{itemize}
\end{frame}

\begin{frame}{Array Implementations: Sorted}
  A priority queue could be implemented as a \emph{sorted} array (again, with a count of elements).
  \pause{}
  \begin{itemize}[<+->]
  \item
    Adding an element would take \(O(n)\) time. (Why?)
  \item
    Removing an element would take \(O(1)\) time. (Why?)
  \item
    Hence, adding and removing an element takes \(O(n)\) time.
  \item
    Again, this is inefficient.
  \end{itemize}
\end{frame}

\begin{frame}{Linked-List Implementations}
  A priority queue could be implemented linked list that is\ldots{}
\begin{itemize}[<+->]
\item
  \emph{Unsorted}:
  \begin{itemize}
  \item
    Adding an element would take \(O(1)\) time (Why?)
  \item
    Removing an element would take \(O(n)\) time (Why?)
  \end{itemize}
\item
  \emph{Sorted}:
  \begin{itemize}
  \item
    Adding an element would take \(O(n)\) time (Why?)
  \item
    Removing an element would take \(O(1)\) time (Why?)
  \end{itemize}
\item
  As with array representations, adding and removing an element takes \(O(n)\) time.
\item
  Again, these are inefficient implementations.
\end{itemize}
\end{frame}

\begin{frame}{Binary Search Tree Implementations}
\begin{itemize}
\item<2->
  A priority queue could be represented as an \\unbalanced \emph{binary search tree}.
  \begin{itemize}
  \item<3->
    Insertion times would range from \alt<-3>{\(O(?)\)}{\(O(\log n)\)} to \alt<-3>{\(O(?)\)}{\(O(n)\)}
  \item<5->
    Removal times would range from \alt<-5>{\(O(?)\)}{\(O(\log n)\)} to \alt<-5>{\(O(?)\)}{\(O(n)\)}
  \end{itemize}
\item<7->
  A priority queue could be represented as a \emph{balanced} binary search tree.%
  \begin{itemize}
  \item<8->
    Insertion and removal could destroy the balance.
  \item<9->
    We need an algorithm to rebalance the binary tree.
  \item<10->
    Good rebalancing algorithms require only \(O(\log n)\) time but are complicated.
  \end{itemize}
\end{itemize}
\end{frame}

\begin{frame}{Heap Implementation}
  \center%
\begin{columns}%
  \column{0.6\textwidth}
    A priority queue can be implemented as a \emph{heap}, which are\ldots{}
    \begin{itemize}
    \item
      Balanced
    \item
      Left-Justified, and have
    \item
      The \emph{Heap Property}: values are least as large as its children.\\
      (Or smaller if smaller values represent higher priorities.)
    \end{itemize}

  \column{0.34\textwidth}
    \centering
    Heap: All nodes have the heap property.
    \begin{forest}
    [12, for tree={inner sep=2pt, minimum size=2em}
      [8
        [4]
        [7]
      ]
      [3]
    ]
    \end{forest}
\end{columns}
\end{frame}

\section{Heaps}
\begin{frame}{Array Representation of a Heap}
\begin{columns}

\column{0.43\textwidth}
\begin{itemize}
\item
  Left child of node \(i\) is at \(2i + 1\); right child is at \(2i + 2\).\\[0.25em]
  %\begin{itemize}
  %\item
    {\small If the result is larger than the last index, there is no such child.}
  %\end{itemize}
\item
  Parent of node \(i\) is \(\frac{1}{2}(i - 1)\), unless \(i\) == \(0\).
\end{itemize}

\column{0.25\textwidth}\centering
  \begin{forest}%
  [9, for tree={inner sep=2pt, minimum height=8mm}
    [7
      [6]
      [0]
    ]
    [5
      [3]
      [, phantom]
    ]
  ]
  \end{forest}

\column{0.20\textwidth}
  \begin{tikzpicture}[node of array/.append style = {minimum height=7mm, minimum width=9mm,}]
  \ArrayLabel{}
  \ArrayListVertical{9, 7, 5, 6, 0, 3}
  \ArrayNodeVertical[none]{}
  \end{tikzpicture}
  % \begin{tabular}{|*{12}{c|}}
  %   \multicolumn{1}{c}{0} & \multicolumn{1}{c}{1} & \multicolumn{1}{c}{2} & \multicolumn{1}{c}{3} & \multicolumn{1}{c}{4} & \multicolumn{1}{c}{5} & \multicolumn{1}{c}{6} & \multicolumn{1}{c}{7} & \multicolumn{1}{c}{8} & \multicolumn{1}{c}{9} & \multicolumn{1}{c}{10} & \multicolumn{1}{c}{11}\\
  %   \hline
  %   \texttt{9} & \texttt{7} & \texttt{5} & \texttt{6} & \texttt{0} & \texttt{3} &   &   &   &   &    &   \\
  %   \hline
  % \end{tabular}


\end{columns}
\end{frame}


\begin{frame}{Adding Element to the Heap}
  \begin{enumerate}
    \item<+->
      Increase \lstinline{lastIndex} and put the new value there.
    \item<+->
      Reheap the newly added element by \emph{sifting up}.
      \begin{itemize}
        \item
          (a.k.a., \emph{up-heap bubbling} or \emph{percolating up}).
        \item
          Sifting up requires \(O(\log n)\) time.
      \end{itemize}
  \end{enumerate}
  \onslide<3->{
  Thus, adding an element takes \(O(\log n)\) time.}
\end{frame}

\begin{frame}{Removing Element from the Heap}
  \begin{enumerate}
    \item<+->
      Remove the element at location 0.
    \item<+->
      Move the element at location \lstinline{lastIndex} to location 0, and decrement \lstinline{lastIndex}.
    \item<+->
      Reheap the new root node (the one now at location 0) by \emph{sifting down}.
      \begin{itemize}
        \item
          (a.k.a., \emph{down-heap bubbling} or \emph{percolating down}).
        \item
          Sifting down requires \(O(\log n)\) time.
    \end{itemize}
  \end{enumerate}
  \onslide<4->{
  Thus, removing an element takes \(O(\log n)\) time.}
\end{frame}

\section{STL Heaps}

\begin{frame}{Using the C++ STL Heap}
\onslide<+->{This is surprising not that bad.}

\begin{itemize}[<+->]
\item
  A heap is a vector that has been ``heapified.''
  \begin{itemize}
  \item
    Use the \lstinline{make_heap()} function to ``heapifiy'' a vector.
  \end{itemize}
\item
  Further pushes/pops should be done via heap functions.
  \begin{itemize}
  \item
    \lstinline{push_back} + \lstinline{push_heap} / \lstinline{pop_heap}.
  \end{itemize}
\end{itemize}

\onslide<+->{Some example code will help.}
\end{frame}



\begin{frame}[fragile, b]{C++ STL Heap: Create}
\begin{lstlisting}[basicstyle=\footnotesize\ttfamily]
std::vector<int> v1 {20, 30, 40, 25, 15}; // Initial a vector
std::make_heap(v1.begin(), v1.end());     // Convert into a heap

std::cout << "The max value is "  << v1.front() << '\n';
\end{lstlisting}
\pause
\begin{columns}[T]
  \column{0.55\textwidth}
\begin{tikzpicture}
  \ArrayLabel{}
  \alt<-2>{
    \ArrayNode{20}
    \ArrayElLabel<-9mm>{v1.begin()}
    \ArrayList{30, 40, 25, 15}
  }{
    \ArrayNode{40}
    \ArrayElLabel<-9mm>{v1.begin()}
    \ArrayList{30, 20, 25, 15}
  }
  \begin{scope}[array index/.append style={text=white!75!CSUblue}]
    \ArrayNode[none]{}
    \ArrayElLabel<-9mm>{v1.end()}
    \ArrayNode[none]{}
  \end{scope}
  \onslide<3>{}
\end{tikzpicture}

\column{0.35\textwidth}
%\onslide<4>{
\begin{forest}
  for tree={inner sep=2pt, minimum size=2em, s sep=5mm}
  [\alt<-2>{20}{40}, name=root
    [30
      [25]
      [15]
    ]
    [\alt<-2>{40}{20}, name=p2
      [50, name=leaf, phantom]{
        %{\draw[<->,dotted, very thick] (leaf) to[bend left] node[midway,left] {sift 1} (p2);}
      }
      [, phantom]
    ]{
      %{\draw[<->,dotted, very thick] (p2) to[bend right] node[midway,right] {sift 2} (root);}
    }
  ]
\end{forest}%
%}%
\\\vspace{10mm}
\end{columns}
\end{frame}

\begin{frame}[fragile, b]{C++ STL Heap: Push}
\begin{lstlisting}[basicstyle=\footnotesize\ttfamily]
v1.push_back(50); // Add a new element to the vector
std::push_heap(v1.begin(), v1.end()); // reorder the elements

std::cout << "The maximum after push is " << v1.front() << '\n';
\end{lstlisting}

\begin{columns}[T]
  \column{0.55\textwidth}
\begin{tikzpicture}
  \ArrayLabel{}
  \alt<-3>{
    \ArrayNode{40}
    \ArrayElLabel<-9mm>{v1.begin()}
    \ArrayNode{30}
    \ArrayNode{20}
    \ArrayList{25, 15}
    \alt<1> {
      \begin{scope}[array index/.append style={text=white!75!CSUblue}]
        \ArrayNode[none]{}
      \end{scope}
    }{
      \ArrayNode{50}
    }
    \only<1>{\ArrayElLabel<-9mm>{v1.end()}}
    \begin{scope}[array index/.append style={text=white!75!CSUblue}]
      \ArrayNode[none]{}
    \end{scope}
    \only<2->{\ArrayElLabel<-9mm>{v1.end()}}
  }{
    \ArrayNode{50}
    \ArrayElLabel<-9mm>{v1.begin()}
    \ArrayList{30, 40, 25, 15, 20}
    \begin{scope}[array index/.append style={text=white!75!CSUblue}]
      \ArrayNode[none]{}
    \end{scope}
    \ArrayElLabel<-9mm>{v1.end()}
  }
  \only<3>{%
    \ArraySwap<sift 1>{el20}{el50}
    \ArraySwap<sift 2>{el40}{el20}
  }%
  \only<4>{%
    \ArraySwap<sift 1>{el40}{el20}
    \ArraySwap<sift 2>{el50}{el40}
  }%
\end{tikzpicture}
\column{0.35\textwidth}
\alt<1>{
  \begin{forest}
    for tree={inner sep=2pt, minimum size=2em, s sep=5mm}
    [40, name=root
      [30
        [25]
        [15]
      ]
      [20,name=p2
        [50, name=leaf, phantom]
        [, phantom]
      ]
    ]
  \end{forest}%
}{
  \begin{forest}
    for tree={inner sep=2pt, minimum size=2em, s sep=5mm}
    [\alt<-3>{40}{50}, name=root
      [30
        [25]
        [15]
      ]
      [\alt<-3>{20}{40},name=p2
        [\alt<-3>{50}{20}, name=leaf]{%
          \onslide<3->{\draw[<->,dotted, very thick] (leaf) to[bend left] node[midway,left] {sift 1} (p2);}%
        }%
        [, phantom]
      ]{%
        \onslide<3->{\draw[<->,dotted, very thick] (p2) to[bend left] node[midway,left] {sift 2} (root);}%
      }
    ]
  \end{forest}%
}\\\vspace{10mm}
\onslide<4>{}
\end{columns}
\end{frame}

\begin{frame}[fragile, b]{C++ STL Heap: Pop}
\begin{lstlisting}[basicstyle=\footnotesize\ttfamily]
std::pop_heap(v1.begin(), v1.end());  // remove the maximum

// Displaying the maximum element of the heap
std::cout << "The maximum after pop is " << v1.front() << '\n';
\end{lstlisting}
\begin{columns}[T]
  \column{0.55\textwidth}
\begin{tikzpicture}
  \ArrayLabel{}
  \alt<-2>{
    \ArrayNode{50}
    \ArrayElLabel<-9mm>{v1.begin()}
    \ArrayList{30, 40, 25, 15, 20}
    \onslide<2>{\ArraySwap{el50}{el20}}%
  }{
    \alt<-5> {
      \ArrayNode{20}
      \ArrayElLabel<-9mm>{v1.begin()}
      \ArrayList{30, 40, 25, 15}
      \onslide<5>{\ArraySwap<sift>{el20}{el40}}%
      \ArrayNode[CSUgold]{50}
      \onslide<3>{\ArraySwap{el20}{el50}}%
    } {
      \ArrayNode{40}
      \ArrayElLabel<-9mm>{v1.begin()}
      \ArrayList{30, 20, 25, 15}
      \onslide<6>{\ArraySwap<sift>{el40}{el20}}%
      \ArrayNode[CSUgold]{50}
    }
  }
  \begin{scope}[array index/.append style={text=white!75!CSUblue}]
    \ArrayNode[none]{}
  \end{scope}
  \ArrayElLabel<-9mm>{v1.end()}
  \onslide<7>{}
\end{tikzpicture}
\column{0.35\textwidth}
\alt<-2>{
  \begin{forest}
    for tree={inner sep=2pt, minimum size=2em, s sep=5mm}
    [50, name=root
      [30
        [25]
        [15]
      ]
      [40
        [20, name=leaf]{%
          \only<2>{\draw[<->,dotted, very thick] (root) to[bend right=45] node[midway,right] {swap} (leaf);}%
        }
        [, phantom]
      ]
    ]
  \end{forest} %
}{
  \alt<-3>{
    \begin{forest}
      for tree={inner sep=2pt, minimum size=2em, s sep=5mm}
      [20, name=root
        [30
          [25]
          [15]
        ]
        [40
          [50, name=leaf, fill=CSUgold]{%
            \only<3>{\draw[<->,dotted, very thick] (root) to[bend right=45] node[midway,right] {swap} (leaf);}%
          }
          [, phantom]
        ]
      ]
    \end{forest}%
  }{
    \alt<-5>{
      \begin{forest}
        for tree={inner sep=2pt, minimum size=2em, s sep=5mm}
        [20, name=root
          [30
            [25]
            [15]
          ]
          [40, name=p2
            [50, phantom]
            [, phantom]
          ]{%
            \onslide<5>{\draw[<->,dotted, very thick] (p2) to[bend left] node[midway,left] {sift} (root);}%
          }
        ]
      \end{forest}%
    }{
      \begin{forest}
        for tree={inner sep=2pt, minimum size=2em, s sep=5mm}
        [40, name=root
          [30
            [25]
            [15]
          ]
          [20, name=p2
            [50, phantom]
            [, phantom]
          ]{%
            \onslide<6>{\draw[<->,dotted, very thick] (p2) to[bend left] node[midway,left] {sift} (root);}%
          }
        ]
      \end{forest}%
    }
  }
}\\\vspace{10mm}
\end{columns}
\end{frame}

\begin{frame}{C++ STL Heap Functions in \texttt{<algorithm>}}
\begin{itemize}[<+->]
\item
  \href{https://en.cppreference.com/w/cpp/algorithm/make_heap}{\lstinline{make_heap()}}: function to ``heapifiy'' a vector.
\item
  \lstinline{push_back()} + \href{https://en.cppreference.com/w/cpp/algorithm/push_heap}{\lstinline{push_heap()}}: push to the heap.
\item
  \href{https://en.cppreference.com/w/cpp/algorithm/pop_heap}{\lstinline{pop_heap()}}: function to remove the first element and reheap.
\item
  \lstinline{front()}: get the top element (first in the heap/queue).
\item
  \href{https://en.cppreference.com/w/cpp/algorithm/make_heap}{\lstinline{sort_heap()}}: perform heapsort (the vector is no longer a heap afterward).
\item
  \href{https://en.cppreference.com/w/cpp/algorithm/is_heap}{\lstinline{is_heap()}}: quick test to see if you are still a heap.
\item
  \href{https://en.cppreference.com/w/cpp/algorithm/is_heap_until}{\lstinline{is_heap_until()}}: returns an iterator of what part of the vector is a heap.

\end{itemize}
\end{frame}

\section{Considerations}
\begin{frame}{Considerations}
\begin{itemize}[<+->]
\item
  A \emph{priority queue} is a data structure that is designed \\to return elements in order of priority.
\item
  Efficiency is usually measured as the sum of the \\time it takes to add and remove elements.
\item
  Simple implementations take \(O(n)\) time.
\item
  A heap implementation takes \(O(\log n)\) time.
\item
  Thus, for any sort of heavy-duty use, a heap implementation is better.
\end{itemize}
\end{frame}

\begin{frame}{My Recommendation}
  \begin{itemize}[<+->]
    \item
      If you need a priority queue, define it as an ADT\@.
    \item
      Start with a simple implementation.
    \item
      Keep the implementation hidden so you can change it later.
    \item
      If appropriate, consider the \href{https://en.cppreference.com/w/cpp/container/priority_queue}{\texttt{std::priority\_queue}} class, which is part of the C++ STL\@.
  \end{itemize}
\end{frame}

\section*{Lab 20}
\begin{frame}{Lab 20: Priority Queues}
\begin{center}\Large
Let's look at the lab based on today's lecture.
\end{center}
\end{frame}
\end{document}
